% 张三慧第4版第6章《振动》独立解析
% V1原始扫描物理页209--212。
% 题目与图纸来源分开记录，专家复核及印刷授权均待完成。
\section*{第6章\quad 振动：独立解析}

\paragraph{习题 6.1}
振动形式 \(x=A\cos(\omega t+\varphi_0)\)，
与 \(x=0.05\cos(8\pi t+\pi/3)\) 比较，
\[
 \boxed{A=0.050\,\mathrm m,\quad
 \omega=8\pi\,\mathrm{rad/s},\quad
 T=\frac{2\pi}{8\pi}=0.250\,\mathrm s,\quad
 \varphi_0=\pi/3}.
\]
由
\(v=-A\omega\sin(\omega t+\varphi_0)\)，
\(a=-A\omega^2\cos(\omega t+\varphi_0)\)，
最大速率与加速度大小
\[
 \boxed{v_{\max}=A\omega=0.4\pi
 \approx1.257\,\mathrm{m/s}},\qquad
 \boxed{a_{\max}=A\omega^2
 =3.2\pi^2\approx31.6\,\mathrm{m/s^2}}.
\]
三个时刻的相位分别为
\[
 \varphi(1)=8\pi+\pi/3,\qquad
 \varphi(2)=16\pi+\pi/3,\qquad
 \varphi(10)=80\pi+\pi/3.
\]
均与 \(\pi/3\) 模 \(2\pi\) 同余。
三条时间曲线可按精确公式
\[
 \boxed{x(t)=0.05\cos(8\pi t+\pi/3)},\quad
 \boxed{v(t)=-0.4\pi\sin(8\pi t+\pi/3)},
\]
\[
 \boxed{a(t)=-3.2\pi^2\cos(8\pi t+\pi/3)}
\]
作图。三图周期相同，速度与位移相差
\(\pi/2\)，加速度与位移反相。
图像的矢量重绘仍待排版。

\paragraph{习题 6.2}
令 \(x=A\cos(\omega t+\varphi_0)\)，
则初速度
\(v_0=-A\omega\sin\varphi_0\)。
(1) \(x_0=-A\)，
必有
\(\boxed{\varphi_0=\pi\pmod{2\pi}}\)。
(2) \(x_0=0,v_0>0\)，
\(\cos\varphi_0=0\) 且
\(\sin\varphi_0<0\)，
故
\(\boxed{\varphi_0=3\pi/2\pmod{2\pi}}\)。
(3) \(x_0=-A/2,v_0<0\)，
\(\cos\varphi_0=-1/2\) 且
\(\sin\varphi_0>0\)，
故
\(\boxed{\varphi_0=2\pi/3\pmod{2\pi}}\)。
相量从正 \(x\) 轴逆时针度量，
它在水平轴上的投影给出 \(x/A\)。

\paragraph{习题 6.3}
原图6.21显示振幅
\(A=5.0\,\mathrm{cm}\)，
\(t=0\) 位移 \(+2.5\,\mathrm{cm}\)
且速度为正。
故初相满足
\(\cos\varphi_0=1/2,\sin\varphi_0<0\)，
取 \(\varphi_0=-\pi/3\)。
图中 \(t=1.0\,\mathrm s\)
为\emph{向下经过} \(x=0\)，
相位为 \(\pi/2\)。
因此
\[
 \omega(1.0)-\pi/3=\pi/2
 \Longrightarrow
 \boxed{\omega=5\pi/6\,\mathrm{rad/s}}.
\]
所以
\[
 \boxed{x(t)=5.0\cos\left(\frac{5\pi t}{6}
 -\frac\pi3\right)\ \mathrm{cm}},\qquad
 \boxed{T=\frac{12}{5}=2.40\,\mathrm s}.
\]
校验图中 \(t=2.2\,\mathrm s\) 时
相位为 \(3\pi/2\)，确为\emph{向上}经过零点，
不能把它误读成 \(x=2.5\,\mathrm{cm}\)。
图中标出的五种状态按循环相位顺序
\[
 \boxed{\varphi_a=0,\quad
 \varphi_b=\pi/3,\quad
 \varphi_c=\pi/2,\quad
 \varphi_d=2\pi/3,\quad
 \varphi_e=4\pi/3}
\]
（均可加 \(2k\pi\)），
相量图按上述角度画出即可。
其中图上各时刻位置仍须随矢量化原图逐点复核。

\paragraph{习题 6.4}
有 \(v_{\max}=A\omega\)，
故角频率
\(\omega=0.03/0.02=1.5\,\mathrm{rad/s}\)，
\[
 \boxed{T=\frac{2\pi}{1.5}
 \approx4.19\,\mathrm s},\qquad
 \boxed{a_{\max}=A\omega^2
 =0.045\,\mathrm{m/s^2}}.
\]
题设 \(t=0\) 为\emph{正向最大速度}，
此时 \(x=0,v>0\)，应取
\(\varphi_0=-\pi/2\)，故
\[
 \boxed{x(t)=0.02\cos(1.5t-\pi/2)
 =0.02\sin(1.5t)\ \mathrm m}.
\]

\paragraph{习题 6.5}
角频率
\(\omega=2\pi/T=4\pi\,\mathrm{rad/s}\)，
振幅 \(A=0.020\,\mathrm m\)。
(1) 初位移
\(x_0=A/2\)，初速度为负。
故初相在第一象限，
\(\boxed{\varphi_{01}=\pi/3}\)：
\[
 \boxed{x_1(t)=0.020\cos(4\pi t+\pi/3)\ \mathrm m}.
\]
(2) 初位移 \(x_0=-A/2\)，
初速度为正。
故初相在第三象限，
\(\boxed{\varphi_{02}=4\pi/3}\)：
\[
 \boxed{x_2(t)=0.020\cos(4\pi t+4\pi/3)\ \mathrm m}.
\]
各自把 \(t=0\) 代入速度式可以验证方向。

\paragraph{习题 6.6}
令第一振子相位
\(\psi_1=\omega t+\varphi\)。
它从正位移返回平衡位置时
\(\psi_1=\pi/2\pmod{2\pi}\)，
第二振子此刻处于正最大位移，
故相位
\(\psi_2=0\pmod{2\pi}\)。
因此第二振子\emph{落后}第一振子
\(\pi/2\)，
\[
 \boxed{x_2=A\cos(\omega t+\varphi-\pi/2)},\qquad
 \boxed{\Delta\varphi=\varphi_2-\varphi_1=-\pi/2}.
\]
第二问若理解「向 \(x\) 方向」为
\emph{正 \(x\) 方向}，那么
\(x_1(0)=-A/2\) 且 \(v_1(0)>0\)，
初相 \(\varphi=4\pi/3\)；
第二振子初相 \(5\pi/6\)。
若原书拟表达负 \(x\) 方向，
则 \(\varphi=2\pi/3\)、
第二振子初相 \(\pi/6\)。
上述两组都保留为可画的条件性相量图，
不可因原句缺少方向符号就擅自确定唯一曲线。

\paragraph{习题 6.7}
取某次相遇时两质点的位置同为
\(x=A/2\)，且速度相反。
同一时刻二者相位可分别取
\(\psi_1=\pi/3\)、
\(\psi_2=-\pi/3\)（或 \(5\pi/3\)）。
于是
\[
 \boxed{|\Delta\varphi|=\frac{2\pi}{3}
 \quad\text{（最小相位差的大小）}}.
\]
若按有向相差规范到 \([0,2\pi)\)，
可能写 \(2\pi/3\) 或 \(4\pi/3\)，
取决于振子的编号及领先/落后约定。
在另一半振幅 \(x=-A/2\) 处也有相同条件。
相量图是两等长相量夹角 \(120^\circ\)。

\paragraph{习题 6.8}
初总能量
\(E=K_0+U_0=0.2+0.6=0.8\,\mathrm J\)
在无阻尼振动中不变。
由 \(E=kA^2/2\) 得
\[
 \boxed{A=\sqrt{\frac{2E}{k}}
 =\sqrt{1.6/25}\approx0.253\,\mathrm m}.
\]
当动、势能相等时
\(U=E/2=0.4\,\mathrm J\)，
\[
 \boxed{|x|=\sqrt{2U/k}
 =\sqrt{0.8/25}\approx0.179\,\mathrm m}.
\]
动、势能之比
\(K/U=(E-\frac12kx^2)/(\frac12kx^2)\)。
在平衡位置 \(x=0\)，势能为零，
此比值没有有限定义、从非零位移趋近时
趋于 \(+\infty\)；如果把题意理解为
\emph{动能最大}，其最大值发生在
\(\boxed{x=0}\)。
不能报告一个凭空的有限「最大比值」。

\paragraph{习题 6.9}
取悬挂弹簧原长位置以下伸长为 \(z\)，
静态平衡满足 \(kb=mg\)。
以平衡位置为原点
\(x=z-b\)，运动方程
\[
 m\ddot x=mg-k(b+x)=-kx,
\quad
 \boxed{\ddot x+\frac{k}{m}x=0}.
\]
因此
\[
 \boxed{\omega=\sqrt{k/m}},\qquad
 \boxed{T=2\pi\sqrt{m/k}}.
\]
由于弹性势能从 \(z=b\) 作差
\(\Delta U_s=\frac12k[(b+x)^2-b^2]
 =kbx+\frac12kx^2\)，
重力势能差 \(\Delta U_g=-mgx\)，
且 \(kb=mg\)，两势能相加正好
\(U=\frac12kx^2\)。
若振幅为 \(A\)，总能量
\[
 \boxed{E=\frac12m\dot x^2+\frac12kx^2
 =\frac12kA^2}.
\]
此能量等式已经用题设的平衡位置势能零点，
而不是在原长位置任意丢弃重力。

\paragraph{习题 6.10}
设弹簧沿其全长的均匀坐标
\(s\in[0,L]\)，单位长度质量 \(m'/L\)，
弹簧下端和振子小球的瞬时速度为 \(v\)。
题设速度线性分布
\(v(s)=vs/L\)。
所以弹簧自身动能
\[
 K_s=\frac12\int_0^L
 \frac{m'}L\left(\frac{vs}L\right)^2ds
 =\frac{m'v^2}{6}.
\]
质点动能 \(\frac12mv^2\)，
合并为
\[
 \boxed{K=\frac12(m+m'/3)v^2}.
\]
因此有效惯性质量
\(m_{\rm eff}=m+m'/3\)。
以小位移的弹性势能
\(U=\frac12kx^2\)，
\[
 \boxed{\omega=\sqrt{\frac{k}{m+m'/3}}}.
\]
要求弹簧各点速度确实服从所给线性分布；
大振幅或弹簧内部波动显著时不再精确。

\paragraph{习题 6.11}
串联两弹簧受力同为 \(F\)，
总伸长
\(x=F/k_1+F/k_2\)，
等效劲度系数
\[
 k_{\rm eq}=\frac{F}{x}
 =\frac{k_1k_2}{k_1+k_2}.
\]
重力只移动平衡位置，
故竖直自由振动周期
\[
 \boxed{T=2\pi\sqrt{\frac m{k_{\rm eq}}}
 =2\pi\sqrt{\frac{m(k_1+k_2)}{k_1k_2}}}.
\]
两根弹簧都按轻质线性弹簧处理。

\paragraph{习题 6.12}
取滑块平衡位置为 \(x=0\)。
沿 \(+x\) 位移 \(x\) 时，
左弹簧伸长带来向左
\(-k_1x\)，右弹簧压缩带来向左
\(-k_2x\)，所以
\[
 m\ddot x=-(k_1+k_2)x,
\quad
 \boxed{\ddot x+\frac{k_1+k_2}{m}x=0}.
\]
两弹簧沿位移方向的恢复力相加，
因此
\[
 \boxed{T=2\pi\sqrt{\frac m{k_1+k_2}}}.
\]
不可套用串联弹簧的等效刚度。

\paragraph{习题 6.13}
两个等质量小球
\(m_1=m_2=m=0.05\,\mathrm{kg}\)，
用质量忽略的弹簧连接，
劲度系数
\(k=1000\,\mathrm{N/m}\)。
两小球在质心系反向振动，
令相对伸长 \(x=x_1-x_2-\ell_0\)，
可得
\[
 \ddot x+k\left(\frac1m+\frac1m\right)x=0.
\]
相对运动约化质量
\(\mu=m/2=0.025\,\mathrm{kg}\)，
所以频率
\[
 \boxed{f=\frac1{2\pi}\sqrt{\frac{k}{\mu}}
 =\frac{1}{2\pi}\sqrt{4.0\times10^4}
 \approx31.8\,\mathrm{Hz}}.
\]
两个球在质心系的位移振幅相同而相位相反，
频率相同。

\paragraph{习题 6.14}
地球均匀密度
\(\rho=3M_E/(4\pi R_E^3)\)；
距地心距离为 \(r\) 的小球
仅受到半径 \(r\) 内均匀物质的引力，
\[
 |\boldsymbol F|=\frac{GmM_Er}{R_E^3}
 \quad\text{朝地心}.
\]
沿任意光滑直隧道中点取坐标 \(x\)，
地心引力沿隧道分量恰为
\[
 \boxed{F_x=-\frac{GM_Em}{R_E^3}\,x}.
\]
故隧道内小球作简谐运动：
\[
 \boxed{\omega=\sqrt{\frac{GM_E}{R_E^3}}},
\qquad
 \boxed{T=2\pi\sqrt{\frac{R_E^3}{GM_E}}}.
\]
近地圆轨道卫星因
\(mv^2/R_E=GM_Em/R_E^2\)，
同样有
\[
 \boxed{T_{\rm orbit}=2\pi\sqrt{R_E^3/(GM_E)}}.
\]
该相等关系只适用于题设\emph{均匀地球、
无阻力}模型，真实地球密度分布非均匀。
